\documentstyle[12pt]{article}

\title{A Polynomial Model for Multi-Valued Logics with a Touch of Algebraic
       Geometry and Computer
       Algebra\footnote{Work partially supported by DGES (Spain).}}

\author{Eugenio Roanes-Lozano
        \thanks{eroanes@eucmos.sim.ucm.es .
        Partially supported by NSA (USA).} \and
        Luis M. Laita  \thanks{laita@fi.upm.es}   \and
        Eugenio Roanes-Mac\'{\i}as \thanks{roanes@eucmos.sim.ucm.es}
        }

\date{}

%******************* Entorno ****************
\newtheorem{entorno}{}[subsection]
\newenvironment{teor}{\begin {entorno} {\bf Theorem.-}}{\end{entorno}}
\newenvironment{prop}{\begin {entorno} {\bf Proposition.-}}{\end{entorno}}
\newenvironment{defi}{\begin {entorno} {\bf Definition.-}}{\end{entorno}}
%\newenvironment{nota}{\begin {entorno} {\bf Note.-}}{\end{entorno}}
\newenvironment{ej}{\begin {entorno} {\bf Example.-}}{\end{entorno}}
\newenvironment{obs}{\begin {entorno} {\bf Remark.-}}{\end{entorno}}
\newenvironment{consec}{\begin {entorno} {\bf Consequence.-}}{\end{entorno}}
%\newenvironment{corol}{\begin {entorno} {\bf Corollary.-}}{\end{entorno}}
\newenvironment{lema}{\begin {entorno} {\bf Lemma.-}}{\end{entorno}}
%\newenvironment{idea}{\begin {entorno} {\bf Idea.-}}{\end{entorno}}
%\newenvironment{inter}{\begin {entorno} {\bf Interpretation.-}}{\end{entorno}}
%\newenvironment{notat}{\begin {entorno} {\bf Notation.-}}{\end{entorno}}
\newenvironment{axis}{\begin {entorno} {\bf Axioms.-}}{\end{entorno}}
\newenvironment{axi}{\begin {entorno} {\bf Axiom.-}}{\end{entorno}}
%\newenvironment{impl}{\begin {entorno} {\bf Implementation.-}}{\end{entorno}}
%*************** Nuevos comandos ************
\newcommand{\NN}{\mbox{{\sl I}}\!\mbox{{\sl N}}}  %La N de reales
\newcommand{\RR}{\mbox{{\sl I}}\!\mbox{{\sl R}}}  %La R de reales
\newcommand{\ZZ}{\mbox{{\sl Z}}\!\!\mbox{{\sl Z}}}  %La Z de enteros
%********************************************

\begin{document}

\maketitle

\vspace{-0.8cm}

\begin{center}
  {\dag \ \S \ Dept. Algebra, Univ. Complutense de Madrid  \\
    \ddag \ Dept. Artificial Intelligence, Univ. Polit\'ecnica de Madrid}
\end{center}

\vspace{0.3cm}

\begin{abstract}
In this paper, a polynomial model (residue class ring) for a given p-valued
propositional Logic (p prime), is constructed. This will allow the study
of logical deductions using Computer Algebra techniques (Gr\"obner Bases).
Also, an interesting interpretation of $\models$ and Kleene's style
$\rightarrow$ and their relation from the point of view of Algebraic Geometry
(in terms of algebraic varieties) will be
given. Only modest requirements about the good behaviour of the Logic will be
assumed. \\
\indent This approach makes it possible to move a step forward and treat
Knowledge Based Systems (KBSs) based on multi-valued Logics.
\end{abstract}
\end{document}
